\begin{claim}{AC.38}{Proved: when scalar Frobenius forces the arithmetic test to be geometric}
Let \(C/k\) be a genus-two curve in characteristic two,
let \(P\in C(k)\), and suppose Jacobian Frobenius is the
scalar \([s]\), where \(s=\varepsilon2^r\),
\(\varepsilon\in\{1,-1\}\), and \(\#k=s^2\).
Let \(\ell\) be an odd prime and put
\[
 f=\operatorname{ord}_\ell(2),\qquad
 m=\operatorname{ord}_\ell(s),\qquad
 d=\frac{f}{\gcd(f,m)}.
\]
Every geometric cyclic degree-\(\ell\) cover has a
\(P\)-pointed descent to \(k\), possibly with nonconstant
deck group, and its first constant-deck field has degree
\(m\). Suppose
\[
                  d=1,
        \quad\text{equivalently}\quad
          \langle2\rangle\subseteq\langle s\rangle
                         \text{ in }\mathbf F_\ell^\times.
\]
Then all its geometric new two-torsion points are rational over
that first constant-deck field. In particular this
cover satisfies (R) over a field with constant deck
group if and only if its geometric new two-rank is zero.
This equivalence holds individually for every cyclic
cover, and persists over all further constant fields.
No assumption that two is primitive modulo \(\ell\)
is required.

The condition has the following explicit forms:
\[
\begin{array}{c|c}
 s&\text{condition for }d=1\\ \hline
 +2^r&\gcd(r,f)=1\\
 -2^r,\ f\text{ even}&\gcd(r+f/2,f)=1\\
 -2^r,\ f\text{ odd}&\gcd(r,f)=1.
\end{array}
\]
When \(d>1\), the same argument yields a norm in
\(\mathbf F_{2^d}^\times\), which need not be one.
It therefore gives no converse obstruction in that case.
\end{claim}
\begin{proof}
The scalar action preserves every geometric degree-\(\ell\)
cyclic quotient of the prime-to-two geometric fundamental
group. The rational point supplies its pointed descent,
with arithmetic Frobenius acting on the deck quotient
by the scalar \(s\), up to the harmless simultaneous
inverse convention. Thus the deck group first becomes
constant after degree \(m\).

Consider the new part of geometric etale cohomology
\[
 W=H^1_{\rm et}(\overline D,\mathbf F_2)_{\rm new},
 \qquad R=\mathbf F_2[X]/\Phi_\ell(X).
\]
The averaging idempotent is defined because \(\ell\)
is odd, so \(W\) is an \(R\)-module.
The Artin--Schreier comparison identifies its scalar
extension with the Frobenius-stable part of the
corresponding coherent cohomology. Every nontrivial
character of the geometric cyclic cover occurs in
\(H^1(\overline D,\mathcal O_{\overline D})\) with
dimension \(g(C)-1=1\). Consequently \(W\) has rank
zero or one over each simple field factor of \(R\).
This is the only step using genus two; compare the
functorial stable-part comparison in AC.17.

Arithmetic Frobenius \(F\) is an automorphism of \(W\)
and is semilinear for
\(\sigma_s:R\to R\), \(X\mapsto X^s\).
If \(W\ne0\), its support is a union of \(\sigma_s\)-orbits
of simple factors. Thus \(W\) is free of rank one over a
\(\sigma_s\)-stable direct factor \(R'\) of \(R\).
Choose a basis on that factor; then
\[
                 F=c\sigma_s\quad(c\in(R')^\times),
 \qquad
                 F^m=\prod_{j=0}^{m-1}\sigma_s^j(c).
\]
The last product belongs to \(((R')^{\sigma_s})^\times\).

Here is the exact fixed-ring calculation. The primitive
\(\ell\)-th roots are indexed by
\(\mathbf F_\ell^\times\); their two-Frobenius orbits
have length \(f\), while multiplication by \(s\) has
order \(m\). As both acting subgroups lie in the same
cyclic group, their intersection has order \(\gcd(f,m)\).
It follows that
\[
 R^{\sigma_s}\simeq
 \left(\mathbf F_{2^{\,f/\gcd(f,m)}}\right)^{
                    (\ell-1)/\operatorname{lcm}(f,m)}.
\]
For \(d=1\) its unit group, and that of every direct
factor such as \((R')^{\sigma_s}\), is trivial.
Therefore \(F^m=1\) on \(W\).
The same holds on the dual geometric new two-torsion
of the Jacobian. If \(W\ne0\), criterion (R) fails at
the first constant field and at every later one. If
\(W=0\), the geometric new two-primary group is already
zero, so (R) holds over every extension. This proves the
equivalence, including its field quantifier.

For \(s=2^r\), one has \(m=f/\gcd(r,f)\), giving
\(d=\gcd(r,f)\). If \(s=-2^r\) and \(f\) is even,
then \(-1=2^{f/2}\pmod\ell\), so replace \(r\) by
\(r+f/2\). If \(f\) is odd, inversion is outside
\(\langle2\rangle\), and
\(m=2f/\gcd(r,f)\); because \(f\) is odd this gives
\(d=\gcd(r,f)\) again.
\end{proof}

\paragraph{Relation to the ordinary-new counterexample.}
For \(C:y^2+y=x^5\) over \(\mathbf F_{16}\), one has
\(s=-4\) and \(r=2\). At \(\ell=7\), \(f=3\) and
\(d=1\); at \(\ell=11\), \(f=10\) and
\(d=\gcd(2+5,10)=1\). For those primes the arithmetic
criterion cannot improve on the geometric new-rank-zero
criterion. At \(\ell=13\), however, \(f=12\) and
\(d=\gcd(2+6,12)=4\). The semilinear norm lies in
\(\mathbf F_{16}^\times\), and AC.31 exhibits the
nontrivial values of order three which give an ordinary-new
rational (R) witness. Thus AC.38 explains exactly why
the unrestricted scalar-Frobenius assertion is false,
while recovering a substantial range in which it is true.
